\section{The complete finite prime-action algebra inside the actual
global
remainder}\label{the-complete-finite-prime-action-algebra-inside-the-actual-global-remainder}

24 September 2026. Complete derivation FPO0--FPO8. The original Riemann
zeta function remains the arithmetic function throughout. The Gaussian
is an explicitly constructed test function; it does not replace zeta or
its spectrum.

\subsection{FPO0. The question answered and the controlling
source}\label{fpo0.-the-question-answered-and-the-controlling-source}

The user's amended instruction USR-64a88219a3ecddd3 is: ``But what
would, based on this finding and everything else we know so far?'' The
preceding PGD calculation constructed a nonzero global dual-comparison
class and each of its positive-real translates. The accumulated
original-zeta trace formula and the user's full-spectrum/dilation
arguments make the next question precise: can a finite linear
combination of distinct global arithmetic rescalings disappear in that
actual quotient? Here we prove that it cannot. We also prove the exact
comparison with the original, unshifted Connes--Consani explicit
formula, including its endpoint contribution.

The source definitions remain B1--B5, their proved consequences P1--P5,
and retraction R1 in
\href{../foundations/user_definitions_20260924/SOURCE_OPERATIONS_AND_PROOFS.md}{\nolinkurl{SOURCE_OPERATIONS_AND_PROOFS.md}}.
The full retained user corpus supplies USR-4322be19bff532cd and the
global arguments USR-9ad1c0a2d09dba92, USR-f55d16feb948d8c2 and
USR-6152e3bc6302258c. These are references into complete preserved
passages, not replacement definitions. Source \(Z_0\) is absence;
\(Z_1/\tau\) is primitive presence with no parity or source addition;
the reconstructed integer layer retains its amounts and \(Z_2\) data. No
operand below is primitive \(\tau\). The arithmetic constant function
\(1\) is not identified with it.

Human sources used here are Alain Connes and Caterina Consani,
\href{https://arxiv.org/abs/2006.13771v1}{\emph{Weil positivity and
Trace formula, the archimedean place}, arXiv:2006.13771v1}, original
author TeX \nolinkurl{weil-compo.tex}, Appendix ``Explicit formula,''
labels \texttt{mellin}, \texttt{bombieriexplicit},
\texttt{bombieriexplicit1} and \texttt{bombieriexplicit2},
lines2033--2048; and their
\href{https://arxiv.org/abs/0903.2024v3}{\emph{Schemes over F1 and zeta
functions}, arXiv:0903.2024v3, §5}, for the receiving coefficient
geometry. The former explicitly credits Enrico Bombieri's \emph{The
Riemann hypothesis} (2006) for its explicit-formula presentation; we
have read that citation in their author source, not claimed a fresh
reading of Bombieri's whole chapter. Ralf Meyer's
\href{https://arxiv.org/abs/math/0412277v3}{\emph{A spectral
interpretation for the zeros of the Riemann zeta function}} supplies the
prior analytic framework. Pierre Deligne's
\href{https://www.numdam.org/item/PMIHES_1980__52__137_0/}{\emph{La
conjecture de Weil II}, §§3.6.1--3.6.3} remains the geometric lifting
target. The new quotient calculation below is our derivation, not a
theorem attributed to those authors.

Complete programme dependencies are
\href{GLOBAL_POLYNOMIAL_DUAL_DEFECT.md}{PGD0--PGD7},
\href{GLOBAL_GAUSSIAN_ORIGINAL_ZETA_TRACE.md}{GGT0--GGT8},
\href{GLOBAL_ORIGINAL_ZETA_RESIDUE_DUALITY_INDEPENDENT.md}{GZR0--GZR9}
and \href{ACTUAL_SUPPORTED_DUALITY_INDEPENDENT.md}{ASD0--ASD14}. Their
full proof bodies remain in the cumulative manuscript.

\subsection{FPO1. Actual domains, ideals and the contour
comparison}\label{fpo1.-actual-domains-ideals-and-the-contour-comparison}

Retain the original spaces, without replacing their topology by an
unrestricted zero-jet product: \[
\mathcal B=\{F\text{ entire}:\sup_{|\sigma|\le A,t\in\mathbb R}
 (1+|t|)^N|F(\sigma+it)|<\infty\text{ for every }A,N\ge0\},
\] \[
\mathcal M=\{H\text{ entire}:\text{for each }A\ge0\text{ some }C_A,d_A
 \text{ satisfy }|H(\sigma+it)|\le C_A(1+|t|)^{d_A}
 \text{ for }|\sigma|\le A\}.
\tag{FPO1.1}
\] Let \(\mathscr Z\) be the original nontrivial zero set, with every
actual multiplicity \(m_\rho\). Set \[
\mathcal J=\{H\in\mathcal M:H^{(j)}(\rho)=0\ (\rho\in\mathscr Z,0\le j<m_\rho)\},
\quad I=\mathcal B\cap\mathcal J,
\quad Q=\mathcal B/I.
\tag{FPO1.2}
\] The quotient \(Q\) carries its proved Fréchet quotient topology; its
continuous complex-linear dual is \(Q'\). Each fixed member of
\(\mathcal M\) multiplies \(\mathcal B\) continuously into itself: on
each strip use its stated polynomial bound and increase the exponent in
the defining seminorm. Thus \(\mathcal B\) is an ideal of
\(\mathcal M\). Leibniz's formula proves that \(\mathcal J\) is an ideal
too, with all multiplicities retained. Consequently \[
\mathcal A_\zeta=\mathcal M/(\mathcal B+\mathcal J)
\tag{FPO1.3}
\] is an algebraic quotient algebra. Its unit is nonzero by PGD4. No
topology on this algebraic quotient is assumed.

The actual continuous-dual contour is \[
\lambda_H([F])=\frac1{2\pi i}
 \left(\int_{2-i\infty}^{2+i\infty}-\int_{-1-i\infty}^{-1+i\infty}\right)
 \frac{H(s)F(1-s)}{\zeta(s)}\,ds.
\tag{FPO1.4}
\] Both edges point upward, and the two integrals converge separately.
PGD1--PGD3 proves descent and continuity on \(Q\), computes
\(\ker(H\mapsto\lambda_H)=\mathcal J\) using all actual zero jets, and
proves the injective linear map \[
\overline\Theta:\mathcal A_\zeta\hookrightarrow
\mathcal C_\zeta=(\chi\otimes Q')/D Q,
\qquad [H]\longmapsto[\lambda_H],\qquad D[F]=\lambda_F.
\tag{FPO1.5}
\] Here the actual action is \[
(T_a^\diamond\lambda)([F])=a\lambda([a^{-s}F(s)]),\qquad a>0,
\quad T_a^\diamond\lambda_H=\lambda_{a^sH}.
\tag{FPO1.6}
\] The factor is \(a\,a^{-(1-s)}=a^s\) in the integrand, so this is a
full group identity, not merely an inference from a formal generator.
The cokernel has been constructed as a vector space with this action; we
do not assert that it carries a multiplication extending that of
\(\mathcal A_\zeta\).

\subsection{FPO2. Construct the original Mellin test and both endpoint
values}\label{fpo2.-construct-the-original-mellin-test-and-both-endpoint-values}

For \(T\ge1\) and \(x>0\), define \[
f_{T,x}(u)=\frac{T}{2\sqrt\pi}
 \exp\!\left(-\frac{T^2}{4}\log^2\frac ux\right),\qquad u>0.
\tag{FPO2.1}
\] Its original Mellin transform, with no shift or completion, is \[
\Phi_{T,x}(s)=\int_0^\infty f_{T,x}(u)u^s\frac{du}{u}
 =x^s e^{s^2/T^2}.
\tag{FPO2.2}
\] For a proof put \(v=\log(u/x)\). The Gaussian integral
\(\int_{\mathbb R}e^{-T^2v^2/4+sv}dv=(2\sqrt\pi/T)e^{s^2/T^2}\) first
follows for real \(s\) by a real translation and completing the square.
Both sides are entire in \(s\); differentiation of the integral is
justified by a Gaussian bound uniformly on compact \(s\)-sets. The
identity theorem gives the stated equality for complex \(s\). On a
bounded vertical strip its modulus is \[
|\Phi_{T,x}(\sigma+it)|=x^\sigma e^{\sigma^2/T^2}e^{-t^2/T^2},
\] which proves membership in \(\mathcal B\) for each fixed \(T,x\).

Connes--Consani's original involution is \[
f_{T,x}^{\sharp}(u)=u^{-1}f_{T,x}(u^{-1})
 =u^{-1}\frac{T}{2\sqrt\pi}
 \exp\!\left(-\frac{T^2}{4}\log^2(ux)\right).
\tag{FPO2.3}
\] Substitution \(v=u^{-1}\) proves its Mellin transform is
\(\Phi_{T,x}(1-s)\). In particular the full two endpoint contributions
are \[
\int_0^\infty f_{T,x}(u)\,du=\Phi_{T,x}(1)=xe^{1/T^2},
\qquad
\int_0^\infty f_{T,x}^{\sharp}(u)\,du=\Phi_{T,x}(0)=1.
\tag{FPO2.4}
\] The latter \(1\) is an integral value. It is different from the point
value \(f_{T,x}(1)=T e^{-T^2\log^2x/4}/(2\sqrt\pi)\); neither is
primitive \(\tau\).

\subsection{FPO3. Prove the domain extension before using the source
formula}\label{fpo3.-prove-the-domain-extension-before-using-the-source-formula}

The cited author source states its explicit formula for
\(C_c^\infty(\mathbb R_+^\times)\). Our \(f_{T,x}\) has noncompact
support. Here is the complete extension needed for this one test family.

Write \(g(v)=f_{T,x}(e^v)\). Fix a smooth function \(\eta\), equal to1
on \([-1,1]\), equal to0 outside \([-2,2]\), and with values in
\([0,1]\). For \(R\ge2\), put \(g_R(v)=\eta(v/R)g(v)\),
\(f_R(u)=g_R(\log u)\). Then \(f_R\) belongs to the source domain. Every
derivative of \(g\) is a polynomial in \(v-\log x\) times its original
Gaussian and its stated constants. It follows directly, by the product
rule and Gaussian decay, that for every \(A\ge0\) and integer \(k\ge0\),
\[
\int_{\mathbb R}e^{A|v|}|(g_R-g)^{(k)}(v)|\,dv\longrightarrow0,
\quad
\sup_v e^{A|v|}|(g_R-g)^{(k)}(v)|\longrightarrow0.
\tag{FPO3.1}
\] Indeed all differences are supported in \(|v|\ge R\); derivatives of
the cutoff contribute powers \(R^{-1}\), uniformly bounded for
\(R\ge2\), and the Gaussian dominates every polynomial and every fixed
exponential on this tail.

For \(|\sigma|\le A\), integrate the Fourier integral for the Mellin
transform by parts \(k\) times in \(v\), placing derivatives on
\(e^{\sigma v}(g_R-g)\). All boundary terms vanish by FPO3.1. Leibniz's
rule bounds the resulting integral by a finite sum of FPO3.1 with
coefficients bounded for \(|\sigma|\le A\). Treat \(|t|\le1\) by the
case \(k=0\). Thus \[
\widetilde f_R\longrightarrow\Phi_{T,x}\quad\text{in }\mathcal B.
\tag{FPO3.2}
\] The actual zero count used in GGT is \(N(Y)\le C(Y+2)^{3/2}\), with
multiplicities. In particular
\(\sum_\rho m_\rho(1+|\Im\rho|)^{-3}<\infty\), as follows by splitting
heights into intervals \([2^j,2^{j+1})\). The strip \(0<\Re\rho<1\) and
FPO3.2 therefore imply convergence of the full zero trace. Both endpoint
evaluations converge too.

Write \(\Lambda(n)=\log p\) when \(n=p^m\) with \(p\) prime and
\(m\ge1\), and0 otherwise. This definition uses the already
reconstructed integer arithmetic and its unique prime factorization. It
has \(0\le\Lambda(n)\le\log n\). The estimate in FPO3.1 with any \(A>2\)
bounds the differences of the first prime series by a constant tending
to0 times \(\sum_{n\ge2}(\log n)n^{-A}\), and bounds the reflected
series by the same type of convergent sum with \(n^{-A-1}\). Hence both
full prime sums converge under this cutoff limit, including every
repetition.

For the original archimedean term, substitute \(u=e^v\) into the
author's exact formula: \[
\mathcal W_{\mathbb R}(f)=
 (\log4\pi+\gamma)g(0)+
 \int_0^\infty
 \frac{g(v)+e^{-v}g(-v)-2e^{-v}g(0)}{1-e^{-2v}}\,dv.
\tag{FPO3.3}
\] The numerator vanishes at \(v=0\), and its derivative is bounded
near0, so its quotient has a removable integrable endpoint bound. For
\(g_R-g\), the numerator is identically0 near0, since \(R\ge2\); its
value at0 is0. Away from0 the denominator is bounded away from0, and
FPO3.1 makes the integral of its difference tend to0. The original tail
term \(-2e^{-v}g(0)\) is retained and integrable. This proves
convergence of the entire archimedean expression.

Every term in the compactly supported source equality thus has its
proved limit. In this equality the zero sum is over nontrivial zeros, as
in the source's subsequent Proposition \texttt{mainprop}. The wording
``all complex zeros'' at source line2042 must not be interpreted as
adding the negative even zeros a second time: the archimedean expression
already incorporates their contribution. GGT3 independently verifies
this convention by the original-zeta contour and every finite leftward
shift. We use no problematic extension of a sum over negative even
integers.

\subsection{FPO4. Exact original prime formula and the retained Gamma
endpoint}\label{fpo4.-exact-original-prime-formula-and-the-retained-gamma-endpoint}

Define the absolutely convergent trace \[
S_T(x)=\sum_{\rho\in\mathscr Z}m_\rho x^\rho e^{\rho^2/T^2}.
\tag{FPO4.1}
\] The source comparison just proved gives, with every coefficient
retained, \[
\begin{aligned}
S_T(x)={}&1+xe^{1/T^2}
 -\frac{T}{2\sqrt\pi}\sum_{n\ge2}\Lambda(n)
 \left\{e^{-T^2\log^2(x/n)/4}
       +n^{-1}e^{-T^2\log^2(xn)/4}\right\}
 -\mathcal W_{\mathbb R}(f_{T,x}).
\end{aligned}
\tag{FPO4.2}
\] The independent calculation GGT1--GGT3 works directly with the
original logarithmic derivative and proves \[
\begin{aligned}
S_T(x)={}&xe^{1/T^2}
 -\frac{T}{2\sqrt\pi}\sum_{n\ge2}\Lambda(n)
 \left\{e^{-T^2\log^2(x/n)/4}
       +n^{-1}e^{-T^2\log^2(xn)/4}\right\}-A_T(x),\\
A_T(x)={}&\frac1{2\pi i}\int_{-1-i\infty}^{-1+i\infty}
 x^s e^{s^2/T^2}\frac{\chi'}{\chi}(s)\,ds,\\
\chi(s)={}&\pi^{s-1/2}\frac{\Gamma((1-s)/2)}{\Gamma(s/2)},\\
\frac{\chi'}{\chi}(s)={}&\log\pi-\tfrac12\psi((1-s)/2)-\tfrac12\psi(s/2).
\end{aligned}
\tag{FPO4.3}
\] Here \(\chi\) is the full functional-equation multiplier; it is not a
replacement arithmetic function. Comparing the two proved equalities
yields the exact relation \[
\boxed{\mathcal W_{\mathbb R}(f_{T,x})=A_T(x)+1.}
\tag{FPO4.4}
\] There is also a direct check: GGT3 proves that shifting the
archimedean contour from \(\Re s=-1\) to \(\Re s=1/2\) crosses exactly
the residue \(\Phi_{T,x}(0)=1\) of \(\chi'/\chi\) at0. Thus
\(A_T=A_T^{(1/2)}-1\), with the stated orientation. No infinite leftward
contour limit has been taken. A finite shift across a trivial zero
\(-2r\) keeps its exact term \(x^{-2r}e^{4r^2/T^2}\) and the remaining
contour, as GGT3 proves.

GGT4--GGT7, from an explicit digamma integral and the same original
formula, establishes the bounds needed below: \[
S_T(1)=\frac{T}{2\sqrt\pi}
 \left(\log T-\log(4\pi)-\frac\gamma2\right)+1+O(\log T),
\tag{FPO4.5}
\] \[
S_T(x)=-\frac{T}{2\sqrt\pi}
 \left\{\Lambda(x)\mathbf1_{x\in\mathbb N,\,x\ge2}
       +x\Lambda(1/x)\mathbf1_{1/x\in\mathbb N,\,1/x\ge2}\right\}
 +x+O_x(\log T)\quad(x\ne1).
\tag{FPO4.6}
\] The symbol \(\Lambda\) in each summand is evaluated only when its
indicated integer condition holds; the whole summand is0 otherwise. All
zero multiplicities are retained, no zero is presumed simple or on the
critical line, and every estimate is for fixed \(x>0\) as
\(T\to\infty\). The distinction used next is exactly the nonzero
coefficient of \(T\log T\) at1 and its absence at every other fixed
positive scale.

\subsection{FPO5. Construct the group algebra and prove its faithful
quotient
map}\label{fpo5.-construct-the-group-algebra-and-prove-its-faithful-quotient-map}

Let \(\mathscr G=\mathbb R_+^\times\), with its usual multiplication.
Define \(\mathbb C[\mathscr G]\) to be the space of functions
\(a:\mathscr G\to\mathbb C\) of finite support. Addition and scalar
multiplication are pointwise; multiplication is finite convolution \[
(a*b)(u)=\sum_{rt=u}a(r)b(t).
\tag{FPO5.1}
\] Only pairs from the two finite supports contribute. Associativity
follows by expanding the same finite triple sum over \(rst=u\); the
function \(\delta_1\) is the multiplicative identity and
\(\delta_r\delta_t=\delta_{rt}\). These operations are constructed in a
receiving arithmetic algebra, not on primitive presence.

For finite support define \[
\mathscr E(a)(s)=\sum_r a(r)r^s.
\tag{FPO5.2}
\] Each summand is entire and bounded on each vertical strip by a
constant depending on that strip and \(r\); hence
\(\mathscr E(a)\in\mathcal M\). Finite expansion proves
\(\mathscr E(a*b)=\mathscr E(a)\mathscr E(b)\) and
\(\mathscr E(\delta_1)=1\). Composition with the ideal quotient gives a
unital algebra map \[
\iota_{\rm alg}:\mathbb C[\mathscr G]\longrightarrow\mathcal A_\zeta.
\tag{FPO5.3}
\] We now prove it is injective; this is stronger than proving the
distinct functions in FPO5.2 are different before quotienting.

Suppose its value at \(a\) is0. List the distinct support points
\(r_1,\ldots,r_k\) and write \(a_i=a(r_i)\). By the definition of
\(\mathcal A_\zeta\), some \(F\in\mathcal B\) satisfies \[
\sum_{i=1}^k a_i r_i^\rho=F(\rho)\qquad(\rho\in\mathscr Z).
\tag{FPO5.4}
\] The quotient requires all multiplicity jets too; the displayed value
equalities are consequences sufficient for the argument. Fix \(j\),
multiply these equalities by \(m_\rho r_j^{-\rho}e^{\rho^2/T^2}\), and
sum over all actual zeros. All these series converge absolutely: each
finite exponential sum has bounded modulus on \(0\le\Re s\le1\), the
Gaussian damps the height, and the zero count above applies. The result
is \[
\sum_{i=1}^k a_i S_T(r_i/r_j)
 =\sum_\rho m_\rho r_j^{-\rho}F(\rho)e^{\rho^2/T^2}.
\tag{FPO5.5}
\] The right side is bounded independently of \(T\ge1\). Indeed \[
|r_j^{-\rho}|\le\max(1,r_j^{-1}),\quad
|e^{\rho^2/T^2}|\le e,
\] and its sum is bounded by \[
e\max(1,r_j^{-1})\,b_{1,3}(F)
 \sum_\rho m_\rho(1+|\Im\rho|)^{-3}<\infty.
\tag{FPO5.6}
\] In the left side, exactly the summand \(i=j\) has scale1, since the
\(r_i\) are distinct. FPO4.5 and FPO4.6 therefore give \[
\sum_{i=1}^k a_i S_T(r_i/r_j)
 =a_j\frac{T\log T}{2\sqrt\pi}+O_a(T).
\tag{FPO5.7}
\] All scales and coefficients are fixed; the finite error constants can
be added. Divide FPO5.5 by \(T\log T\) and let \(T\to\infty\). Its right
side tends to0 and its left side tends to \(a_j/(2\sqrt\pi)\). Therefore
\(a_j=0\). Repeat for each \(j\); hence \(a=0\). This proves \[
\boxed{\mathbb C[\mathbb R_+^\times]\hookrightarrow\mathcal A_\zeta
 \hookrightarrow\mathcal C_\zeta,\qquad
 \delta_r\longmapsto[r^s]\longmapsto c_r=[\lambda_{r^s}].}
\tag{FPO5.8}
\] The first arrow is an injective algebra map. The second and their
composition are injective linear maps; no multiplication on all of
\(\mathcal C_\zeta\) has been asserted.

\subsection{FPO6. The entire prime algebra and every finite
relation}\label{fpo6.-the-entire-prime-algebra-and-every-finite-relation}

The subgroup \(\mathbb Q_+^\times\subset\mathbb R_+^\times\) consists
exactly of finite products \(\prod_p p^{n_p}\), \(n_p\in\mathbb Z\),
uniquely. To check uniqueness, write a proposed equality as an equality
of two positive integers by multiplying away the finite negative
exponents, and apply the retained unique integer prime factorization.
Conversely every positive rational has such a presentation by factoring
its numerator and denominator.

It follows, with the displayed multiplication rather than a naming
convention, that \[
\mathbb C[\mathbb Q_+^\times]
 \simeq\mathbb C[X_p,X_p^{-1}:p\text{ prime}],
\qquad \prod_p X_p^{n_p}\longmapsto
 \delta_{\prod_p p^{n_p}}.
\tag{FPO6.1}
\] Both sides consist of finite sums of finite monomials. Uniqueness of
the exponents proves the map is bijective, and addition of exponents
proves it preserves multiplication. Its exact injection into the actual
cokernel is \[
\prod_p X_p^{n_p}\longmapsto
 \left[\lambda_{\exp(s\sum_p n_p\log p)}\right].
\tag{FPO6.2}
\] Every logarithm and exponent has its original positive-rational
value. No common rescaling of the arithmetic times is introduced.

FPO1.6 proves that \(T_a^\diamond c_r=c_{ar}\). Thus the injected space
is the regular group module: multiplication by \(\delta_a\) corresponds
to the actual positive-real action. Its orbit of \(c_1=c_\zeta\) is
exactly the basis \(\{c_r:r>0\}\). In particular, for every prime \(p\),
all nonzero Laurent polynomials in \(T_p^\diamond\) act nontrivially on
\(c_\zeta\). For distinct primes, the same holds for every nonzero
multivariate Laurent polynomial, using the same proved
prime-factorization bijection. This statement concerns their action on
this class, not invertibility of those operators on the whole cokernel.

Moreover \(T_a^\diamond c_\zeta=c_\zeta\) forces \(a=1\), since
otherwise \(\delta_a-\delta_1\ne0\) would map to0. This is an exact
statement about detectability in this named receiver. It does not
declare that every geometrical change of presentation of the user's
global counting construction is this particular action, or that all such
presentations have been classified.

\subsection{FPO7. Exact scope of the supported comparison and source
labels}\label{fpo7.-exact-scope-of-the-supported-comparison-and-source-labels}

ASD14 proves the complete cone of the original coefficient comparison
has degree-one cohomology \(\mathcal C_\zeta\). On representatives the
final map into that cohomology is \[
\delta_r\longmapsto\lambda_{r^s}\longmapsto\pi'\lambda_{r^s}
 \longmapsto[\pi'\lambda_{r^s}]\in H^1(K_\zeta),
\tag{FPO7.1}
\] where \(\pi:A\to Q\) is the actual raw Mellin quotient. The
representative is a cycle since \(\pi d=0\) gives \(d'\pi'\lambda=0\). A
boundary would mean \(\pi'\lambda=\pi'D\pi(b)\); injectivity of \(\pi'\)
would imply \(\lambda\in DQ\), exactly the relation excluded above. This
proves both the actual map and its injectivity in the cone, not only in
a newly defined algebra.

The source receiving-ring action is retained on the full coefficient
complex, with its original extra closed copies. On this spectral term
its action is the first arithmetic scalar action, hence it commutes with
every displayed linear map. That fact does not identify source \(\tau\)
with integer1: their different actions on the faithful extra closed
copies remain part of the full complex and are not quotiented away here.

For the existing receiving support lattice \(L\), the exact lift of any
displayed linear map \(f:V\to W\) is \[
(v,\alpha)\longmapsto(f(v),\alpha),\qquad
G_L(V)=\{(0,\alpha):\alpha\in L\}\cup(V\times\{1_L\}).
\tag{FPO7.2}
\] A nontop input has zero amplitude, which linearity sends to zero; a
nonzero output therefore has top label. Substitution proves preservation
of identities and composition, and injectivity is retained when \(f\) is
injective. A labelled zero is never turned into an unlabelled absence.
This is the named receiving construction, not a newly imposed operation
on \(Z_1/\tau\).

\subsection{FPO8. Contribution to the current lifting
calculation}\label{fpo8.-contribution-to-the-current-lifting-calculation}

PGD constructed individual global orbit elements. This proof determines
their full finite relation space: it is zero. It constructs the entire
real group algebra and the Laurent algebra of all primes as an injected
subalgebra of the original multiplier quotient, and their exact
equivariant linear injection into the actual dual-comparison cone. The
test derives from the original prime formula and retains both endpoint
values, the complete Gamma contribution and every prime-power
repetition. The exact relation \(\mathcal W_{\mathbb R}=A_T+1\) prevents
losing the original endpoint when moving between the two presentations.

The numerical weight separation sought in the active tau goal is not
asserted by this representation theorem. The next calculation is the
actual support-localization map on these now distinguished global
classes, together with the other degrees of the full complex: it must
determine which map would have to annihilate an obstruction to a lifted
class. That calculation uses the already constructed ASD localization
row and complete source-faithful complex, rather than treating the whole
nonzero comparison cone as a single finite-weight obstruction. The
complete
\href{GLOBAL_UNIT_ORBIT_AND_SUPPORTED_COMPARISON.md}{UOS3--UOS8A
derivation} carries it out: the exact cone map sends this entire module
by \(c_r\mapsto(c_r,0)\) into
\(H^1(\mathcal L_\zeta)=\mathcal C_\zeta\oplus(\chi\otimes Q')\); it is
the identity on \(H^{-1}=H\), and is the explicit restriction
\(\chi H'\to\chi E'\) in degree2. The original Fourier restriction has
its written section and zero connecting map. Thus these relative
residue-comparison classes have a proved supported placement, and have
not been substituted for a failure of that existing Fourier lift.
